\chapter{Fixed Points and Conformal Data}
\label{ch:volume-iii-fixed-points-conformal-data}

\section{Local QFT Data Used in Fixed-Point Analysis}

Let \(D\geq2\) be the spacetime dimension, let \(d=D-1\), and let
\(\eta_{\mu\nu}=\operatorname{diag}(-,+,\ldots,+)\).  The local Lorentzian
data used in this volume are the following.
\begin{enumerate}
  \item A Hilbert space \(\Hilb\) of physical states.
  \item A strongly continuous unitary representation \(U(\Lambda,a)\) of the
  proper orthochronous Poincare group, or of its spin cover when spinorial
  fields occur.  The group law is
  \[
    U(\Lambda',a')U(\Lambda,a)
    =
    U(\Lambda'\Lambda,\Lambda'a+a') .
  \]
  \item Self-adjoint generators \(P_\mu\) and \(J_{\mu\nu}=-J_{\nu\mu}\),
  normalized near the identity by
  \[
    U(\Lambda,a)
    =
    1+\ii a^\mu P_\mu
    +\frac{\ii}{2}\omega^{\mu\nu}J_{\mu\nu}
    +O(a^2,\omega^2,a\omega),
    \qquad
    \Lambda^\mu{}_\nu=\delta^\mu{}_\nu+\omega^\mu{}_\nu .
  \]
  The operators \(P_\mu\) generate energy and momentum, while \(J_{\mu\nu}\)
  generates spatial rotations and Lorentz boosts.
  \item A Poincare-invariant unit vector \(\vac\in\Hilb\), so
  \(U(\Lambda,a)\vac=\vac\).
  \item Local fields \(\widehat\Phi_\alpha(x)\), understood as
  operator-valued distributions on a common invariant dense domain, with
  Poincare covariance
  \[
    U(\Lambda,a)^{-1}\widehat\Phi_\alpha(x)U(\Lambda,a)
    =
    R(\Lambda)_\alpha{}^\beta\,
    \widehat\Phi_\beta(\Lambda x+a).
  \]
  Here \(R\) is the finite-dimensional Lorentz representation carried by the
  field component label \(\alpha\).  Bosonic local fields commute, and
  fermionic local fields anticommute, after smearing with test functions whose
  supports are spacelike separated.
\end{enumerate}

In addition, the discussion of conformal symmetry assumes a renormalized
local stress tensor \(\widehat T_{\mu\nu}(x)\).  It is a symmetric tensor-valued
operator distribution satisfying
\begin{equation}
  \partial_\mu\widehat T^{\mu\nu}=0,
  \qquad
  \widehat T^{\mu\nu}=\widehat T^{\nu\mu}.
  \label{eq:cft-opening-stress-conservation}
\end{equation}
where the conservation equation is distributional.  When the spatial
integrals exist with the chosen domain and infrared prescription, the
translation and Lorentz generators are represented by
\begin{equation}
  P^\mu
  =
  \int_{\Sigma_t}\dd^d\vec x\,
  \widehat T^{0\mu}(t,\vec x).
  \label{eq:cft-opening-p-charge}
\end{equation}
and
\begin{equation}
  J^{\mu\nu}
  =
  \int_{\Sigma_t}\dd^d\vec x\,
  \bigl(
    x^\mu\widehat T^{0\nu}(t,\vec x)
    -
    x^\nu\widehat T^{0\mu}(t,\vec x)
  \bigr).
  \label{eq:cft-opening-j-charge}
\end{equation}
with \(\Sigma_t=\{x^0=t\}\).  These equations are operator equations after
the same smearing, renormalization, and boundary conventions used to define
\(\widehat T_{\mu\nu}\).

Lorentz covariance of the stress tensor means
\begin{equation}
  U(\Lambda,0)\widehat T_{\mu\nu}(x)U(\Lambda,0)^{-1}
  =
  \Lambda_\mu{}^\rho\Lambda_\nu{}^\sigma
  \widehat T_{\rho\sigma}(\Lambda x).
  \label{eq:cft-opening-stress-lorentz}
\end{equation}
where indices on \(\Lambda\) are moved with \(\eta_{\mu\nu}\).  Hence
\(\widehat T^\mu{}_\mu(x)\) is a scalar local operator:
\[
  U(\Lambda,0)\widehat T^\mu{}_\mu(x)U(\Lambda,0)^{-1}
  =
  \widehat T^\mu{}_\mu(\Lambda x).
\]

\section{Renormalized Operators and the Trace Equation}

Fix a renormalized coordinate chart \(\{g_I(\mu)\}_{I\in\mathcal I}\) on a
space of local scalar deformations.  The index set \(\mathcal I\) labels
renormalized local insertions \([O_I]_\mu(x)\).  In a Lagrangian presentation
with bare scalar operators \(O_I^{\rm bare}\), write the interaction density as
\begin{equation}
  \mathcal L_{\rm int}^{\rm bare}
  =
  \sum_I g_I^{\rm bare}O_I^{\rm bare}.
  \label{eq:cft-opening-bare-interaction}
\end{equation}
The renormalized insertions are defined by the response of the bare density to
renormalized coordinates:
\begin{equation}
  \delta\mathcal L_{\rm int}^{\rm bare}
  =
  \sum_I \delta g_I^{\rm bare}O_I^{\rm bare}
  =
  \sum_I \delta g_I(\mu)\,[O_I]_\mu .
  \label{eq:cft-opening-renormalized-insertions}
\end{equation}
This equation fixes the normalization and mixing convention for the local
operators in the chosen scheme.  The beta-function vector field is
\begin{equation}
  \frac{\dd}{\dd\log\mu}g_I(\mu)
  =
  \beta_I(g(\mu)).
  \label{eq:cft-opening-beta-function}
\end{equation}

The trace equation is the local form of the response to a change of scale.
For the MS-type source charts developed earlier, the derivation has two
steps.  First, the regulated scale variation of the bare interaction produces
the engineering-dimension factors multiplying \(g_I^{\rm bare}O_I^{\rm bare}\).
Second, the identity
\[
  0
  =
  \frac{\dd g_I^{\rm bare}}{\dd\log\mu}
  =
  -(\text{engineering dimension of }g_I^{\rm bare})\,g_I^{\rm bare}
  +
  \sum_J\beta_J(g)
  \frac{\partial g_I^{\rm bare}}{\partial g_J}
  \]
converts those engineering factors into the beta-function vector field acting
on the same bare density.  Passing to the finite renormalized insertion
defined in \eqref{eq:cft-opening-renormalized-insertions} gives, in flat space
and away from coincident insertions,
\begin{equation}
  \widehat T^\mu{}_\mu
  =
  \sum_I\beta_I(g)\,[O_I]_\mu
  +
  \partial_\mu V^\mu .
  \label{eq:cft-opening-trace-beta}
\end{equation}
The vector \(V^\mu\) is a local vector operator recording possible virial and
improvement contributions in the chosen stress-tensor representative.  Contact
terms in \eqref{eq:cft-opening-trace-beta} are fixed together with the
renormalized source Ward identities.  Curvature terms appear after coupling
the theory to a background metric, but vanish in the flat separated-point
identity displayed here.

An RG fixed point in the coordinate chart is a point \(g_\ast\) satisfying
\begin{equation}
  \beta_I(g_\ast)=0
  \qquad
  \text{for every }I\in\mathcal I .
  \label{eq:cft-opening-fixed-point}
\end{equation}
At such a point the beta-function part of the trace vanishes.  A flat-space
conformal fixed point is obtained when the stress tensor can be chosen, by
allowed improvements and contact-term conventions, so that
\begin{equation}
  \widehat T^\mu{}_\mu=0
  \label{eq:cft-opening-traceless-stress}
\end{equation}
as a renormalized local operator equation in flat space.  Equation
\eqref{eq:cft-opening-traceless-stress}, together with
\eqref{eq:cft-opening-stress-conservation}, is the local input from which the
conformal currents and charges are derived in the next chapter.

The linearized RG flow near \(g_\ast\) records the scaling dimensions of
deformations.  Let
\[
  B_{IJ}
  =
  \left.
  \frac{\partial\beta_I}{\partial g_J}
  \right|_{g=g_\ast}.
\]
Choose an eigenvector \(v^{(a)}_J\) and write
\[
  g_I-g_{\ast I}=\lambda_a v^{(a)}_I+O(\lambda^2),
  \qquad
  \sum_J B_{IJ}v^{(a)}_J=-y_a v^{(a)}_I .
\]
With the momentum scale \(\mu\) as flow parameter,
\[
  \mu\frac{\dd\lambda_a}{\dd\mu}
  =
  -y_a\lambda_a+O(\lambda^2).
\]
Equivalently, under a length rescaling \(x\mapsto b x\), \(b>1\),
\[
  \lambda_a(b)=b^{y_a}\lambda_a+O(\lambda^2).
\]
If the corresponding scalar perturbation is
\(\int\dd^Dx\,\lambda_a O_a(x)\), then
\begin{equation}
  y_a=D-\Delta_a,
  \label{eq:cft-opening-y-delta}
\end{equation}
where \(\Delta_a\) is the scaling dimension of \(O_a\).  Thus \(y_a>0\),
\(y_a=0\), and \(y_a<0\) define relevant, marginal, and irrelevant
directions, respectively, in the long-distance flow.

\section{Conformal Field Theory as Local Fixed-Point Data}

The local data of a conformal field theory in flat \(D\)-dimensional space are
organized as follows.
\begin{enumerate}
  \item A vector space \(\mathcal V_{\rm loc}\) of renormalized local
  operators.  Its elements are denoted
  \(\mathcal O_A(x)\), with \(A\) labeling a basis or a collection of
  multiplets.
  \item A stress tensor \(T_{\mu\nu}\) satisfying conservation, symmetry, and
  the flat-space tracelessness condition
  \[
    \partial_\mu T^{\mu\nu}=0,
    \qquad
    T^{\mu\nu}=T^{\nu\mu},
    \qquad
    T^\mu{}_\mu=0 .
  \]
  \item An action of the conformal algebra on \(\mathcal V_{\rm loc}\),
  generated by the stress-tensor charges associated with conformal Killing
  fields.
  \item Euclidean separated-point correlation functions
  \[
    \left\langle
    \mathcal O_{A_1}(x_1)\cdots\mathcal O_{A_n}(x_n)
    \right\rangle,
    \qquad
    x_i\in\R^D,
  \]
  interpreted as distributions on the configuration space
  \(x_i\neq x_j\), together with contact-term extensions specified by the
  source Ward identities.
\end{enumerate}
For unitary Lorentzian theories these Euclidean distributions are related to
Wightman distributions by analytic continuation and reflection positivity,
under the same reconstruction hypotheses used for ordinary Euclidean Green
functions.

\section{Source Functionals and Ward-Identity Contact Terms}
\label{sec:cft-source-functionals-ward-contacts}

The separated-point correlators in the preceding list are restrictions of
distributions obtained from a source functional.  Work in Euclidean signature
for this construction.  The fixed-point functional used here is the
renormalized metric-and-operator source limit constructed in
Section~\ref{sec:background-metric-source-functional}, specialized to a point
of the RG flow where the beta functions vanish and to a stress-tensor
improvement convention appropriate to the fixed point.  Thus \(g_{\mu\nu}\) is
a smooth background metric, \(J^A(x)\) are smooth compactly supported
sources for a chosen local operator basis \(\mathcal O_A\), and local
functionals of \(g\) and \(J\) specify the contact-term chart.  The sign
normalization of \(W_\ast\) is the Euclidean free-energy, or effective-action,
normalization \(W_\ast=-\log Z_\ast\).  The connected-log functional is denoted
\(K_\ast=\log Z_\ast\) when it is used.
For a regulated path-integral representative this convention reads
\[
  \exp\{-W_\ast[g,J]\}
  =
  \int \mathcal D_g\Phi\,
  \exp\left\{
    -S_\ast[\Phi;g]
    +\int\dd^D x\,\sqrt g\,J^A(x)\mathcal O_A(x)
  \right\},
\]
up to the local counterterms that define the source chart.  The minus sign in
the operator derivative below is forced by combining \(W_\ast=-\log Z_\ast\)
with the source term \(+\int J^A\mathcal O_A\) in the exponent; a source written
with the opposite sign would reverse this derivative convention.  In an abstract
CFT the displayed integral is replaced by the corresponding renormalized source
functional; the first derivatives below are the definition of the associated
contact-term convention.
The stress tensor and local operators are the source derivatives
\begin{equation}
  \left\langle T^{\mu\nu}(x)\right\rangle_{g,J}
  =
  -\frac{2}{\sqrt g(x)}
  \frac{\delta W_\ast[g,J]}{\delta g_{\mu\nu}(x)},
  \qquad
  \left\langle \mathcal O_A(x)\right\rangle_{g,J}
  =
  -\frac{1}{\sqrt g(x)}
  \frac{\delta W_\ast[g,J]}{\delta J^A(x)} .
  \label{eq:cft-source-derivative-conventions}
\end{equation}
The stress-tensor derivative is written with respect to \(g_{\mu\nu}\).  The
same convention may be written using \(g^{\mu\nu}\) by the chain-rule identity
\eqref{eq:stress-tensor-covariant-contravariant-derivatives}.
Higher derivatives define Euclidean Schwinger distributions.  The source
functional is defined up to local functionals of the background sources.  This
is the continuum CFT version of the contact-term coordinate freedom described in
Section~\ref{sec:multiple-insertions-contact-coordinates}: the restriction to
the configuration space \(x_i\ne x_j\) is fixed by the linear operator
normalization, while the extension across collision diagonals is fixed by the
chosen local source chart.

The Ward identities are identities of these source-dependent distributions.
For an infinitesimal diffeomorphism generated by a compactly supported vector
field \(\xi^\mu\),
\[
  \delta_\xi g_{\mu\nu}
  =
  \nabla_\mu\xi_\nu+\nabla_\nu\xi_\mu,
  \qquad
  \delta_\xi J^A
  =
  \xi^\mu\nabla_\mu J^A
\]
for scalar sources.  Invariance of \(W_\ast[g,J]\) gives
\begin{equation}
  \nabla_\mu
  \left\langle T^{\mu\nu}(x)\right\rangle_{g,J}
  =
  \left\langle \mathcal O_A(x)\right\rangle_{g,J}
  \nabla^\nu J^A(x),
  \label{eq:cft-source-diffeomorphism-ward}
\end{equation}
with extra representation terms when the source carries spin indices.
Differentiating this identity with respect to \(J\) and then setting
\(J=0\) produces the contact terms in stress-tensor Ward identities.
Equivalently, for scalar insertions
\[
  X=\mathcal O_{A_1}(x_1)\cdots\mathcal O_{A_n}(x_n),
\]
the distributional translation Ward identity can be written in smeared form as
\begin{equation}
  \int\dd^D x\,
  \partial_\mu\xi_\nu(x)\,
  \left\langle T^{\mu\nu}(x)X\right\rangle
  =
  \sum_{i=1}^n
  \left\langle
    \mathcal O_{A_1}(x_1)\cdots
    (\xi^\rho\partial_\rho\mathcal O_{A_i})(x_i)\cdots
    \mathcal O_{A_n}(x_n)
  \right\rangle .
  \label{eq:cft-smeared-translation-ward}
\end{equation}
This equation is a compact way of recording the delta-function terms
supported at \(x=x_i\).

If the fixed point has no flat-space Weyl anomaly, the source chart for scalar
primary operators of dimensions \(\Delta_A\) may be chosen so that an
infinitesimal Weyl transformation satisfies
\[
  \delta_\sigma g_{\mu\nu}=2\sigma g_{\mu\nu},
  \qquad
  \delta_\sigma J^A=(\Delta_A-D)\sigma J^A .
\]
Differentiating the corresponding source Ward identity gives, for scalar
primaries and in flat space,
\begin{equation}
  \int\dd^D x\,\sigma(x)\,
  \left\langle T^\mu{}_\mu(x)X\right\rangle
  =
  -\sum_{i=1}^n
  \Delta_{A_i}\sigma(x_i)
  \left\langle X\right\rangle
  +
  \mathcal A_\sigma[X].
  \label{eq:cft-smeared-weyl-ward}
\end{equation}
The term \(\mathcal A_\sigma[X]\) is a local functional of the sources and
background metric.  It vanishes for flat separated-point correlators in a
Weyl-invariant source scheme, but it must be specified when correlators are
extended to coincident insertions or when the theory is placed on a curved
background.

A scalar primary operator \(\mathcal O\) has a scaling dimension
\(\Delta_{\mathcal O}\).  For scalar primaries
\(\mathcal O_i\), separated-point scale covariance is
\begin{equation}
  \left\langle
  \mathcal O_1(\lambda x_1)\cdots\mathcal O_n(\lambda x_n)
  \right\rangle
  =
  \lambda^{-\sum_i\Delta_i}
  \left\langle
  \mathcal O_1(x_1)\cdots\mathcal O_n(x_n)
  \right\rangle,
  \qquad
  \lambda>0 .
  \label{eq:cft-opening-scale-covariance}
\end{equation}
Operators with spin carry the corresponding finite-dimensional rotation or
Lorentz matrices.  The subsequent chapters derive these transformation laws
from stress-tensor charges.

\section{Ising Universality as a Fixed-Point Example}

The statistical Ising model is a Euclidean lattice model whose microscopic
state sum has a finite-dimensional diagonal trace representation.  Let
\(\Lambda\subset\mathbb Z^D\) be a finite lattice.  At every site
\(x\in\Lambda\) there is a spin variable \(s_x=\pm1\), represented in diagonal
operator notation by
\[
  V_x=\operatorname{span}\{\ket{\uparrow}_x,\ket{\downarrow}_x\},
  \qquad
  \widetilde{\Hilb}_\Lambda=\bigotimes_{x\in\Lambda}V_x,
\]
and by commuting diagonal operators \(\widehat S_x\) with eigenvalues
\(\pm1\).  With nearest-neighbor coupling normalized to one,
\[
  \widetilde H
  =
  -\sum_{\langle x x'\rangle}\widehat S_x\widehat S_{x'} .
\]
The thermal expectation of an observable \(\widehat O\) is
\begin{equation}
  \langle\widehat O\rangle_{\beta_{\rm T}}
  =
  \frac1Z
  \operatorname{Tr}_{\widetilde{\Hilb}_\Lambda}
  \left(
    \ee^{-\beta_{\rm T}\widetilde H}\widehat O
  \right),
  \qquad
  Z=
  \operatorname{Tr}_{\widetilde{\Hilb}_\Lambda}
  \ee^{-\beta_{\rm T}\widetilde H},
  \qquad
  \beta_{\rm T}=\frac1T .
  \label{eq:cft-opening-ising-thermal-trace}
\end{equation}
The tilde distinguishes the finite-dimensional trace space of the statistical
ensemble from the Hilbert space reconstructed from a continuum scaling limit.

\begin{figure}[H]
  \centering
  \resizebox{0.98\textwidth}{!}{%
  \begin{tikzpicture}[x=1cm,y=1cm,every node/.style={font=\scriptsize}]
    \begin{scope}
      \node[font=\small] at (1.35,3.0) {lattice spins};
      \foreach \i in {0,...,4} {
        \draw[gray!55] (0,\i*0.45)--(2.7,\i*0.45);
        \draw[gray!55] (\i*0.675,0)--(\i*0.675,1.8);
      }
      \foreach \x/\y/\s in {
        0/0/1,.675/0/-1,1.35/0/1,2.025/0/1,2.7/0/-1,
        0/.45/-1,.675/.45/1,1.35/.45/1,2.025/.45/-1,2.7/.45/1,
        0/.9/1,.675/.9/1,1.35/.9/-1,2.025/.9/1,2.7/.9/-1,
        0/1.35/-1,.675/1.35/1,1.35/1.35/1,2.025/1.35/1,2.7/1.35/-1,
        0/1.8/1,.675/1.8/-1,1.35/1.8/1,2.025/1.8/-1,2.7/1.8/1} {
        \ifnum\s=1
          \draw[green!50!black,->,thick] (\x,\y-0.15)--(\x,\y+0.15);
        \else
          \draw[green!50!black,->,thick] (\x,\y+0.15)--(\x,\y-0.15);
        \fi
      }
      \node[align=center] at (1.35,-0.45)
        {\(x=(i,j)\) in \(D=2\)\\\(s_x=\pm1\)};
    \end{scope}

    \begin{scope}[shift={(4.45,0)}]
      \node[font=\small] at (1.75,3.0) {magnetization};
      \draw[-{Latex[length=2mm]}, thick] (0,0)--(3.55,0) node[below] {\(\beta_{\rm T}\)};
      \draw[-{Latex[length=2mm]}, thick] (0,0)--(0,2.3) node[left] {\(|\langle s\rangle|\)};
      \draw[dashed,gray!70] (1.23,0)--(1.23,2.05);
      \node[below] at (1.23,0) {\(\beta_c\)};
      \draw[blue!75!black,very thick]
        (0.05,0.02)--(1.23,0.02)
        .. controls (1.30,0.38) and (1.65,1.30) .. (2.45,1.78)
        .. controls (2.85,2.00) and (3.2,2.05) .. (3.45,2.05);
      \node[blue!75!black,align=center] at (2.65,1.15)
        {ordered\\phase};
      \node[gray!70!black] at (-0.18,2.05) {\(1\)};
    \end{scope}

    \begin{scope}[shift={(9.45,0)}]
      \node[font=\small] at (1.75,3.0) {spin two-point function};
      \draw[-{Latex[length=2mm]}, thick] (0,0)--(3.8,0) node[below] {\(r\)};
      \draw[-{Latex[length=2mm]}, thick] (0,0)--(0,2.3) node[left] {\(f(r)\)};
      \draw[green!50!black, very thick]
        (0.15,2.02)..controls (0.85,1.72) and (1.75,1.48)..(3.55,1.46);
      \draw[blue!75!black, very thick]
        (0.15,2.02)..controls (0.55,1.28) and (1.25,0.70)..(3.35,0.28);
      \draw[purple!75!black, very thick]
        (0.15,2.02)..controls (0.45,0.95) and (0.95,0.28)..(2.25,0.05);
      \node[green!50!black] at (2.85,1.78) {\(T<T_c\)};
      \node[blue!75!black] at (2.65,0.73) {\(T=T_c\)};
      \node[purple!75!black] at (1.45,0.43) {\(T>T_c\)};
    \end{scope}
  \end{tikzpicture}
  }
  \caption{The Ising fixed point is reached by tuning the temperature.  The
  finite lattice trace space defines the statistical ensemble, spontaneous
  magnetization appears for \(T<T_c\), and the critical two-point function has
  power-law decay.}
  \label{fig:cft-opening-ising-source-figures}
\end{figure}

In \(D=2\), write \(S_{N,M}\) for the spin observable at
\(x=(N,M)\) and set
\[
  f_{N,M}
  =
  \langle S_{0,0}S_{N,M}\rangle_{\beta_{\rm T}}
  \simeq f(r),
  \qquad
  r=(N^2+M^2)^{1/2},
\]
for \(r\gg1\) and \(T\) close to \(T_c\).  The long-distance regimes are
\[
  f(r)\longrightarrow |\langle s\rangle|^2>0
  \qquad (T<T_c),
\]
\[
  f(r)\sim \ee^{-r/\xi(T)}
  \qquad (T>T_c),
  \qquad
  \xi(T)\sim (T-T_c)^{-\nu},
\]
and
\begin{equation}
  f(r)\sim r^{-2\Delta_\sigma}
  \qquad (T=T_c).
  \label{eq:cft-opening-ising-critical-power}
\end{equation}
The first numerical values to keep in view are
\[
\begin{array}{c|cc}
  D & \Delta_\sigma & \nu \\ \hline
  2 & \frac18 & 1\\[2pt]
  3 & 0.5181488\ldots & 0.629971\ldots
\end{array}
\]
where \(\Delta_\sigma\) is the scaling dimension of the continuum spin field
and \(\nu\) is the correlation-length exponent.
\begin{remark}[Numerical provenance in three dimensions]
The \(D=3\) entries are numerical conformal-bootstrap data, not results derived
in this chapter.  A recent mixed-correlator computation including the stress
tensor gives
\(\Delta_\sigma=0.518148806(24)\) and
\(\Delta_\varepsilon=1.41262528(29)\) for the leading
\(\mathbb Z_2\)-odd and \(\mathbb Z_2\)-even scalar primaries of the critical
Ising CFT; using \(\Delta_\varepsilon=D-1/\nu\) at \(D=3\) gives
\(\nu=0.629970975\ldots\).  The source is Chang, Dommes, Erramilli, Homrich,
Kravchuk, Liu, Mitchell, Poland, and Simmons-Duffin,
\emph{Bootstrapping the 3d Ising Stress Tensor}, arXiv:2411.15300.
\end{remark}

The same universality class can be represented by a scalar Euclidean lattice
field theory.  Replace the two-valued spin by a real variable \(\phi_x\), use
a \(\mathbb Z_2\)-even double-well single-site potential, and retain the
finite Brillouin-zone cutoff.  At long distance its action has the form
\begin{equation}
  S_{\rm lat}[\phi]
  =
  \int\dd^Dx\,
  \left[
    \frac12 Z_\phi(\partial_\mu\phi)^2
    +
    V(\phi)
    +
    \sum_A c_A\,\mathcal O_A^{\rm aniso}(\phi,\partial\phi)
  \right],
  \label{eq:cft-opening-ising-scalar-cutoff}
\end{equation}
where the last sum consists of higher-derivative operators allowed by the
hypercubic lattice symmetry.  These terms display the finite cutoff and the
breaking of \(SO(D)\) rotations at the microscopic scale.

At \(T=T_c\), equivalently at the tuned massless critical point, the
correlation length is infinite in lattice units.  The separated long-distance
correlators are described by the Ising conformal field theory.  For
\(D=2,3\), this is the infrared fixed point of the \(\mathbb Z_2\)-invariant
massless scalar \(\phi^4\) universality class.  The continuum fields obey the
dictionary
\[
  \phi \longmapsto \sigma,
  \qquad
  [\phi^2]\longmapsto \varepsilon ,
\]
where \(\sigma\) is the spin field and \(\varepsilon\) is the thermal
energy-density field after identity mixing has been subtracted.  The thermal
deformation
\[
  \Delta S_t=\int\dd^Dx\,u_t\,\varepsilon(x)
\]
has RG eigenvalue \(y_t=D-\Delta_\varepsilon\).  Since
\(\xi^{-1}\propto |u_t|^{1/y_t}\), comparison with
\(\xi\propto |T-T_c|^{-\nu}\) gives
\begin{equation}
  \Delta_\varepsilon
  =
  D-\frac1{\nu}.
  \label{eq:cft-opening-ising-energy-dimension}
\end{equation}

\begin{figure}[H]
  \centering
  \resizebox{0.96\textwidth}{!}{%
  \begin{tikzpicture}[x=1cm,y=1cm,every node/.style={font=\scriptsize}]
    \begin{scope}
      \node[font=\small] at (0,2.65) {cutoff scalar model};
      \draw[-{Latex[length=2mm]}, thick] (-1.35,0)--(1.45,0) node[below] {\(\phi\)};
      \draw[-{Latex[length=2mm]}, thick] (0,-0.2)--(0,2.05) node[left] {\(V(\phi)\)};
      \draw[blue!70!black,very thick]
        plot[smooth] coordinates {(-1.25,1.85) (-0.72,0.35) (0,1.05) (0.72,0.35) (1.25,1.85)};
      \node[blue!70!black] at (0,-0.55) {\(\mathbb Z_2\)-even};
    \end{scope}

    \begin{scope}[shift={(4.1,0)}]
      \node[font=\small] at (0,2.65) {finite momentum cutoff};
      \draw[thick] (-1.15,-1.15) rectangle (1.15,1.15);
      \draw[-{Latex[length=2mm]}, gray!70] (-1.55,0)--(1.65,0) node[below] {\(k^1\)};
      \draw[-{Latex[length=2mm]}, gray!70] (0,-1.55)--(0,1.65) node[left] {\(k^2\)};
      \node at (-1.15,-1.42) {\(-\pi\)};
      \node at (1.15,-1.42) {\(\pi\)};
      \node[align=center] at (0,-1.9) {Brillouin zone};
    \end{scope}

    \draw[-{Latex[length=2mm]}, thick] (5.75,0.55)--(7.05,0.55);
    \node[align=center] at (6.40,1.05) {tune\\mass};

    \begin{scope}[shift={(8.05,0)}]
      \node[font=\small] at (0,2.65) {infrared fixed point};
      \draw[thick,blue!60!black] (-1.2,0.15)..controls (-0.55,1.7)
        and (0.55,1.7)..(1.2,0.15);
      \draw[thick,blue!60!black,dashed] (-1.2,0.15)..controls (-0.55,-0.38)
        and (0.55,-0.38)..(1.2,0.15);
      \fill[purple!75!black] (-0.42,0.75) circle (2pt);
      \node[purple!75!black,above] at (-0.42,0.86) {\(\sigma\)};
      \fill[orange!85!black] (0.55,0.58) circle (2pt);
      \node[orange!85!black,above] at (0.55,0.69) {\(\varepsilon\)};
      \node[align=center] at (0,-0.70)
        {\(\Delta_\sigma\), \(\Delta_\varepsilon=D-1/\nu\)};
    \end{scope}
  \end{tikzpicture}
  }
  \caption{The scalar cutoff representation of the Ising universality class
  has a finite Brillouin-zone cutoff and lattice anisotropy.  At the tuned
  critical point its long-distance operator data are those of the Ising
  conformal field theory.}
  \label{fig:cft-opening-ising-scalar-cutoff}
\end{figure}
